%% start of file `template.tex'.
%% Copyright 2006-2013 Xavier Danaux (xdanaux@gmail.com).
%
% This work may be distributed and/or modified under the
% conditions of the LaTeX Project Public License version 1.3c,
% available at http://www.latex-project.org/lppl/.


\documentclass[10pt,a4paper,sans]{moderncv}        % possible options include font size ('10pt', '11pt' and '12pt'), paper size ('a4paper', 'letterpaper', 'a5paper', 'legalpaper', 'executivepaper' and 'landscape') and font family ('sans' and 'roman')

% modern themes
\moderncvstyle{banking}                            % style options are 'casual' (default), 'classic', 'oldstyle' and 'banking'
\moderncvcolor{blue}                                % color options 'blue' (default), 'orange', 'green', 'red', 'purple', 'grey' and 'black'
%\renewcommand{\familydefault}{\sfdefault}         % to set the default font; use '\sfdefault' for the default sans serif font, '\rmdefault' for the default roman one, or any tex font name
%\nopagenumbers{}                                  % uncomment to suppress automatic page numbering for CVs longer than one page

% character encoding
\usepackage[utf8]{inputenc}                       % if you are not using xelatex ou lualatex, replace by the encoding you are using
%\usepackage{CJKutf8}                              % if you need to use CJK to typeset your resume in Chinese, Japanese or Korean

% adjust the page margins
\usepackage[scale=0.75]{geometry}
\usepackage{xcolor}
%\setlength{\hintscolumnwidth}{3cm}                % if you want to change the width of the column with the dates
%\setlength{\makecvtitlenamewidth}{10cm}           % for the 'classic' style, if you want to force the width allocated to your name and avoid line breaks. be careful though, the length is normally calculated to avoid any overlap with your personal info; use this at your own typographical risks...

\usepackage{import}
\usepackage{fontawesome}



% personal data
\name{\LARGE{Piotr}}{Maciąg}
%\title{Curriculum Vitae}                               % optional, remove / comment the line if not wanted
%\address{my address, line 1, line 2, line 3, postcode}{}{}% optional, remove / comment the line if not wanted; the "postcode city" and and "country" arguments can be omitted or provided empty
\vspace{1.5pt}
% \phone[mobile]{+48 692 096 723}                   % optional, remove / comment the line if not wanted
%\phone[fixed]{01234 123456}                    % optional, remove / comment the line if not wanted
%\phone[fax]{+3~(456)~789~012}                      % optional, remove / comment the line if not wanted
\email{piotr.maciag@pw.edu.pl}                               % optional, remove / comment the line if not wanted
% \homepage{piotr-maciag.github.io}                         % optional, remove / comment the line if not wanted

\extrainfo{\faLinkedin~linkedin.com/in/piotrmaciag32~
\faGraduationCap~scholar.google.pl/citations?user=v--wmw0AAAAJ}% optional, remove / comment the line if not wanted


%\photo[64pt][0.4pt]{picture}                       % optional, remove / comment the line if not wanted; '64pt' is the height the picture must be resized to, 0.4pt is the thickness of the frame around it (put it to 0pt for no frame) and 'picture' is the name of the picture file
%\quote{Some quote}                                 % optional, remove / comment the line if not wanted

% to show numerical labels in the bibliography (default is to show no labels); only useful if you make citations in your resume
%\makeatletter
%\renewcommand*{\bibliographyitemlabel}{\@biblabel{\arabic{enumiv}}}
%\makeatother
%\renewcommand*{\bibliographyitemlabel}{[\arabic{enumiv}]}% CONSIDER REPLACING THE ABOVE BY THIS

% bibliography with mutiple entries
%\usepackage{multibib}
%\newcites{book,misc}{{Books},{Others}}
%----------------------------------------------------------------------------------
%            content
%----------------------------------------------------------------------------------
\begin{document}
%\begin{CJK*}{UTF8}{gbsn}                          % to typeset your resume in Chinese using CJK
%-----       resume       ---------------------------------------------------------
\makecvtitle

\vspace{-1cm}


\section{Academic Experience}

%\vspace{6pt}
\cventry{08.2024 - now}{Institute of Computer Science, Warsaw University of Technology}{Head of Postgraduate Studies in \textcolor{color1}{AI-Powered Business Intelligence}}{Warsaw, Poland}{}{
 Various administrative duties related to managing studies, lectures, classes and contacts with students.
}
\cventry{10.2022 - now}{Institute of Computer Science, Warsaw University of Technology}{Assistant Professor}{Warsaw, Poland}{}{
    \begin{itemize}
        \item Conduct research in data science topics such as \textcolor{color1}{prediction}, \textcolor{color1}{clustering}, \textcolor{color1}{NLP}, and \textcolor{color1}{pattern discovery}. 
        \item Lecture in \textcolor{color1}{Databases, Artificial Intelligence Methods}, and \textcolor{color1}{Object-Oriented Programming}.
        \item Supervising bachelor and master theses in Computer Science.
    \end{itemize}
}

% \item{\cventry{10.2022 - *}{Assistant Professor}{}{}{}{
%     \begin{itemize}
%         \item Conducting research (preparing scientific publications, research proposals) in various data science topics: prediction, clustering, NLP, patterns discovery.
        
%         \item Supervising undergraduate and graduate students.
%     \end{itemize}
% }}

\cventry{03.2019 - 09.2022}{Institute of Computer Science, Warsaw University of Technology}{Research and Teaching Assistant}{Warsaw, Poland}{}{}

\section{Research Visits}

\cventry{01.2019--02.2019}{Robotics Institute}{Carnegie Mellon University}{Pittsburgh, USA}{}{}

\cventry{01.2019--02.2019}{Optimization, Modeling and Data Analytics Research Group}{Tecnalia Research and Innovation}{Bilbao, Spain}{}{}

\cventry{09.2017--07.2018}{Knowledge Engineering and Discovery Research Institute}{Auckland University of Technology}{Auckland, New Zealand}{}{}




%\item{\cventry{07.2015--09.2015}{Databases and Java programmer}{Comarch}{Rzeszow, Poland}{}{\vspace{3pt} }}



% \section{Working experience}

% \begin{itemize}
%     \item Research and teaching assistant at the Institute of Computer Science, Warsaw University of Technology, Warsaw, Poland; \textit{03.2019 - Present}.
%     \item Visiting researcher at Robotics Institute, Carnegie Mellon University, Pittsburgh, USA; \textit{12.10.2019 - 13.12.2019}.
%         \item Visiting researcher at the OPTIMA research group, Tecnalia, Bilbao, Spain; \textit{20.01.2019 - 16.02.2019}.
%     \item Visiting researcher at Knowledge Engineering and Discovery Research Institute, Auckland University of Technology, New Zealand; \textit{23.09.2017 - 14.07.2018}. 
%     \item Junior Java programmer, Comarch, Rzeszów, Poland; \textit{07.2015 - 08.2015}.
% \end{itemize}

\section{Industry Experience}
\cventry{11.2025 - *}{Processifier}{Process Mining Engineer}{Warsaw, Poland}{}{
Developing process mining applications and dashboards with the use of Celonis platform.
}



\cventry{07.2023 - 12.2023}{CDeX}{AI Specialist}{Warsaw, Poland}{}{
    \begin{itemize}
        \item Developed an AI agent to simulate cybersecurity threats within the CDeX environment.
        \item Technologies I worked with included \textcolor{color1}{Python, MLFlow, XML databases, Docker, Git}.
    \end{itemize}
}

\cventry{08.2022--10.2023}{Accenture}{Data Science Analyst}{Warsaw, Poland}{}{
    \begin{itemize}
        \item Developed and maintained ETL pipelines in a large production environment.
        \item Wrote and optimized SQL queries within the constraints of the client’s cloud-based database environment, including \textcolor{color1}{Hive, MySQL, Presto}, and \textcolor{color1}{Spark} managed in \textcolor{color1}{Mercurial}
        \item Created Business Intelligence dashboards and reports using \textcolor{color1}{Tableau} and \textcolor{color1}{Unidash}. 
    \end{itemize}
}

\section{Education}

\vspace{5pt}

%\subsection{Academic Qualifications}



\cventry{2015--2021}{Ph.D. in Computer Science}{Warsaw University of Technology}{}{\textit{Distinguished PhD Thesis}}{}


\cventry{2014--2015}{M.Sc in Computer Science}{Rzeszów University of Technology}{}{Graduated as a honour student}{}  % arguments 3 to 6 can be left empty

\cventry{2010--2014}{B.Sc in Computer Science}{Rzeszów University of Technology}{}{}{}

\newpage

  \section{Main Technical Skills}
\begin{itemize}
 \item Programming \& Scripting: Python (scikit-learn, PyTorch, Keras), R, SQL, PL/SQL, LaTeX, C++ (in the past).
 \item Machine Learning \& Data Science: Time series forecasting, anomaly detection, clustering, NLP, data mining, complexity of algorithms. 
\item Big Data \& Pipelines: Spark, Hive, Presto, Dataswarm, MySQL, PostgreSQL.
 \item Visualization \& Reporting: Tableau, Unidash, Jupyter Notebooks.
 \item Version Control \& Tools: Git, Mercurial, Docker (basic).
 \end{itemize}

% \vspace{6pt}

% \begin{itemize}

% \item \textbf{Data science/machine learning/databases skills:} Proficient in various classification and regression models and algorithms (SVM, neural networks, decision trees), spiking neural networks, time series processing, anomaly detection, clustering, NLP (in text classification domain), optimization and administration in databases (Oracle, PostgreSQL).

% \vspace{6pt}

% \item \textbf{Proficiency in programming languages and technologies:} SQL, PL/SQL, R (including data visualization packages and data mining/machine learning packages), C++, GIT (GitHub, GitLab, Bitbucket), \LaTeX, Matlab, experience in machine learning with Python (Kears, PyTorch, SNNTorch, scikit-learn), Jupyter notebooks, distributed database environments: Hive, Presto, Spark, tools for orchestration of data engineering pipelines: Dataswarm. Data visualization tools (Tableau).

% \vspace{6pt}

% \item \textbf{General business skills:} several years of experience in teaching technical and computer science classes and lectures for undergraduate and postgraduate students, proficiency in technical writing in English and Polish, many presentations at conferences/seminars/technical talks. 

% \end{itemize}

% \vspace{6pt}



%\section{Projects}
%\vspace{6pt}

% \newpage 

\section{Recent Awards, Scholarships and Projects}
% \vspace{6pt}
 \begin{itemize}
 \item Principal investigator in the project ``Building reusable representations in deep machine learning'' - \textit{Polish National Science Center Sonata funding scheme}, 2025.
 \item Awarded four research grants by the \textit{Excellence Initiative Research University} at Warsaw University of Technology  for the research in data mining, machine learning, neural networks (total awarded funding approx. 300 000 PLN), 2024.
 \item Rector's Award in Science, Warsaw University of Technology, 2021.
 \item Rector's Scholarship for Young Researchers, Warsaw University of Technology, 2021.
\item Research grant awarded by the council of Technical Computer Science and Telecommunications discipline, Warsaw University of Technology, 2021.
 \item Research grant for young researchers awarded by the Faculty of Electronics and Information Technology, Warsaw University of Technology, 2016.
 \end{itemize}

 \section{Main Publications}

% \vspace{6pt}

\begin{enumerate}

    \item \textbf{Piotr S. Maci{\k{a}}g}, ``Mining targeted spatio-temporal sequential patterns'', \textit{GeoInformatica}, pp. 1–29, 2025.

    \item \textbf{Piotr S. Maci\k{a}g}, Robert Bembenik, Aleksandra Piekarzewicz, Javier Del Ser, Jesus L. Lobo, and Nikola K. Kasabov,
    ``Effective air pollution prediction by combining time series decomposition with stacking and bagging ensembles of evolving spiking neural networks'',
    \textit{Environmental Modelling \& Software}, vol. 170, pp. 105851, 2023.

    \item \textbf{Piotr S. Maci\k{a}g}, Robert Bembenik, and Artur Dubrawski,
    ``Discovery of crime event sequences with constricted spatio-temporal sequential patterns'', \textit{Journal of Big Data}, vol. 10, no. 1, pp. 98, 2023.
    
    \item \textbf{Piotr S. Maci\k{a}g}, Nikola Kasabov, Marzena Kryszkiewicz and Robert Bembenik,``Air pollution prediction with clustering-based ensemble of evolving spiking neural networks and a case study for London area'', \textit{Environmental Modelling \& Software}, vol. 118, pp. 262-280, 2019, DOI: 10.1016/j.envsoft.2019.04.012.

    \item \textbf{Piotr S. Maci\k{a}g}, Marzena Kryszkiewicz, Robert Bembenik, Jesus L. Lobo and Javier Del Ser, ``Unsupervised Anomaly Detection in Stream Data with Online Evolving Spiking Neural Networks'', \textit{Neural Networks}, vol. 139, pp. 118–139. 2021.
    
    \item \textbf{Piotr S. Maci\k{a}g}, Marzena Kryszkiewicz and Robert Bembenik, ``Online Evolving Spiking Neural Networks for Incremental Air Pollution Prediction'', in \textit{Proceedings of the 2020 International Joint Conference on Neural Networks: IJCNN 2020}, pp. 1 - 8.
    
    \item\textbf{Piotr S. Maci\k{a}g}, Marzena Kryszkiewicz and Robert Bembenik, ''Discovery of closed spatio-temporal sequential patterns from event data'', in \textit{Proceedings of the 23rd International Conference Knowledge-Based and Intelligent Information {\&} Engineering Systems: KES 2019}, pp. 707-716.
    
     % \item \textbf{Piotr S. Maci\k{a}g}, ''A Survey on Data Mining Methods for Clustering Complex Spatiotemporal Data'', in \textit{Proceedings of 13th International Conference Beyond Databases, Architectures and Structures. Towards Efficient Solutions for Data Analysis and Knowledge Representation: BDAS 2017}, vol. 716, pp. 115–126, Springer International Publishing.

     \item \textbf{Piotr S. Maci\k{a}g}, ''Efficient Discovery of Sequential Patterns from Event-Based Spatio-Temporal Data by Applying Microclustering Approach'', in \textit{Intelligent Methods and Big Data in Industrial Applications}, vol. 40, pp. 183–199, 2019, Springer International Publishing, DOI: 10.1007/978-3-319-77604-0\_14.
    
     % \item \textbf{Piotr S. Maci\k{a}g}, ''Efficient Discovery of Top-K Sequential Patterns in Event-Based Spatio-Temporal Data'', in \textit{Proceedings of the 2018 Federated Conference on Computer Science and Information Systems: FedCSIS 2018}, vol. 15, pp. 47–56, DOI: 10.15439/2018F19.

     % \item \textbf{Piotr S. Maci\k{a}g} and Robert Bembenik, ''A Novel Breadth-first Strategy Algorithm for Discovering Sequential Patterns from Spatio-temporal Data'', in \textit{Proceedings of the 8th International Conference on Pattern Recognition Applications and Methods: ICPRAM 2019}, pp. 459–466, DOI: 10.5220/0007355804590466.


    
     % \item \textbf{Piotr S. Maci\k{a}g}, ''A Framework for Discovering Frequent Event Graphs from Uncertain Event-based Spatio-temporal Data'', in \textit{Proceedings of the 8th International Conference on Pattern Recognition Applications and Methods: ICPRAM 2019}, pp. 656-663, DOI: 10.5220/0007411206560663.
    
\end{enumerate}

You can find my related research publications in the DBLP database: \url{https://dblp.org/pid/200/0320.html} and Google Scholar \url{https://scholar.google.pl/citations?user=v--wmw0AAAAJ&hl=pl}.

 % \section{Courses and Certificates}

\section{Main Presentations at Research Conferences}
\vspace{6pt}

\begin{itemize}
\item Keynote speech ,,Spatio-temporal pattern discovery: Is a data mining approach really enough?" at conference 6th International Conference on Frontiers in Computing and Systems (COMSYS), Warsaw, Poland 25-27 September 2025. 
\item Speech ,,A Comparative Study of Short Text Classification with Spiking Neural Networks ‘ conference 2022 17th Conference on Computer Science and Intelligence Systems (FedCSIS), Sofia, Bulgaria, 04 – 07 September 2022.
\item Speech ,,New theoretical results in design of evolving spiking neural networks and their application to air pollution prediction and anomaly detection’’ congress IEEE World Congress on Computational Intelligence (WCCI) 2020, Glasgow, Wielka Brytania, 19 - 24 lipca 2020.
\item Poster presentation ,,Online Evolving Spiking Neural Networks for Incremental Air Pollution Prediction’’ congress IEEE World Congress on Computational Intelligence (WCCI) 2020, Glasgow, Great Britain, 19 - 24 July 2020.
\item Speech ,,Discovery of closed spatio-temporal sequential patterns from event data’’ conference 23rd International Conference on Knowledge-Based and Intelligent Information \& Engineering Systems (KES 2019), Budapest, Hungary, 4 - 6 September 2019.
\end{itemize}

\section{Other Recent Activities}
\vspace{6pt}

\begin{itemize}
    \item Peer-reviewing publications for journals and conferences: \textit{Artificial Intelligence Review} journal, \textit{IEEE Transactions on Neural Networks and Learning Systems}, \textit{Applied Mathematics and Computation}, \textit{34th ACM/SIGAPP Symposium On Applied Computing (SAC 2019)}, \textit{4th IEEE International Conference on Data Science and Advanced Analytics (DSAA 2017)} and others.

    \item Member of the Program Committee of the scientific workshops “Design, Implementation and Applications of Spiking Neural Networks and Neuromorphic Systems”, organized during the IEEE World Congress on Computational Intelligence 2020, Glasgow, United Kingdom, July 19–24, 2020.
\end{itemize}



% Publications from a BibTeX file without multibib
%  for numerical labels: \renewcommand{\bibliographyitemlabel}{\@biblabel{\arabic{enumiv}}}% CONSIDER MERGING WITH PREAMBLE PART
%  to redefine the heading string ("Publications"): \renewcommand{\refname}{Articles}
\nocite{*}
\bibliographystyle{plain}
\bibliography{publications}                        % 'publications' is the name of a BibTeX file

% Publications from a BibTeX file using the multibib package
%\section{Publications}
%\nocitebook{book1,book2}
%\bibliographystylebook{plain}
%\bibliographybook{publications}                   % 'publications' is the name of a BibTeX file
%\nocitemisc{misc1,misc2,misc3}
%\bibliographystylemisc{plain}
%\bibliographymisc{publications}                   % 'publications' is the name of a BibTeX file

%-----       letter       ---------------------------------------------------------

\end{document}


%% end of file `template.tex'.
